\section{Details of the School Reconstruction Plan}\label{sec:reconstructionplan}
This section provides an English translation of the ``Education" section from a document authored by the Presidential Delegation for Reconstruction, the Ministry of the Interior and Public Security \citep{Gobierno2010}, detailing the reconstruction plan following the 2010 Maule earthquake.\par 

In Table \ref{tab:plan}, I highlight in bold font all mentions of school types targeted by the reconstruction plan. As evident from the table, the plan specifically targeted schools based on the extent of damage at the establishment level.


\subsection*{Translation of governmental policy document}

In educational matters, the earthquake and tsunami of February 27 meant that 2,095,671 students saw their schools damaged, delaying the start of their school year. The disaster-affected area had 8,326, of which 6,168 suffered some kind of damage, corresponding to 74 percent. Forty-eight percent of schools in the affected areas had moderate, severe, or disabling damage.

\begin{table}[H]
\centering
\caption{Summary of Schools Affected by the Earthquake by Region}
\begin{tabular}{|l|c|c|}
\hline
\textbf{Region} & \textbf{Number of Schools} & \textbf{Enrollment} \\
\hline
Valparaíso & 997 & 289,724 \\
O'Higgins & 620 & 166,153 \\
Maule & 732 & 175,469 \\
Biobío & 1,155 & 355,186 \\
Araucanía & 423 & 97,056 \\
Metropolitan & 2,241 & 1,012,082 \\
\hline
\textbf{Total} & 6,168 & 2,095,670 \\
\hline
\end{tabular}
\end{table}

To solve the problems caused by the earthquake, the Ministry of Education developed an Emergency and Reconstruction Plan organized in four stages: emergency, stabilization, early reconstruction, and reconstruction, described as follows:

\begin{table}[h!]
\centering
\caption{Stage, Deadlines, Description}\label{tab:plan}
\begin{tabular}{|l|l|p{8cm}|}
\hline
\textbf{Stage} & \textbf{Deadlines} & \textbf{Description} \\
\hline
Emergency & 27-02-10 to 26-04-10 & Infrastructure habilitation for class commencement. Within 45 days all students returned to classes. \\
\hline
Stabilization & 26-04-2010 to 26-07-11 & Period of reconstruction aimed to stabilize the school system, focusing on reassignment or relocation of temporary facilities used during the emergency period that could affect students' health, did not adequately permit curricular activities (non-educational establishments), or cases where coexistence of school communities sharing infrastructure posed critical situations. Approximately 200 {\bf schools} were detected {\bf in critical conditions}, assigning regional executives who supported local urgent needs according to established criteria, and facilitating appropriate solutions (modular classrooms, insulation of temporary housing adapted during the emergency, among others). \\
\hline
Early Reconstruction & 27-02-10 to 27-02-11 & Stage aimed at normalizing infrastructure for the maximum number of students and supporting municipal and subsidized private schools to recover habitability and safety conditions. For this, the Minor Repairs Plan was launched in July, with around one thousand schools applying and preparing a second stage for late August, totaling 30 billion pesos. \\
\hline
Reconstruction & 27-02-10 to 26-02-14 & The last stage aims to finalize repairs of minor and moderate damage and focus efforts on {\bf severely damaged schools} and replacements, to finish the process before the 2014 school year begins. This Plan, benefiting {\bf schools with severe damage}, both municipal and subsidized private, will launch in September for 30 billion pesos. Additionally, 15 emblematic high schools will benefit this year, with an estimated 35 billion pesos, for infrastructure repair projects. They will be selected through joint work between municipalities, Regional Ministerial Secretaries of Education, and the community, using criteria of high social and local recognition and at least one thousand students. \\
\hline
\end{tabular}
\end{table}
